

\chapter{Introduction}\label{cap:intro}

\begin{quote}
    If physical theories were people, thermodynamics would be the village witch. [...] The other theories find her somewhat odd, somehow different in nature from the rest, yet everyone comes to her for advice, and no-one dares to contradict her.\\
    - Goold, Huber, Riera, del Rio, Skrzypczyk (2016)
\end{quote}

The reason everyone comes to this ``village witch" for advise it is her pragmatic approach to nature. Instead of trying to describe microscopic details of a system, she focus on what can be done with its macroscopic resources. This is a consequence of thermodynamics being developed in the context of thermal machines and the goal to achieve efficient exploitation of the natural world. Moreover, this theory development brought fundamental laws regarding energy flow and transformation that have been used in different areas of physics, chemistry, biology, and other sciences. Despite thermodynamics describing macroscopic systems near equilibrium, there is a great effort to extend its concepts to nonequilibrium and quantum regime \cite{goold-2016, vinjanampathy-2016, parrondo-2015, deffner-2019}. This connection of quantum mechanics and thermodynamics can help to understand how quantum systems exchange energy with its surroundings. Because quantum properties can affect transient dynamics, relaxation processes of quantum systems become a non-trivial and timely problem.

Relaxation processes are studied in different areas of physics. Some usual examples are resistor-capacitor circuit, damped harmonic oscillators, dielectric relaxation, and thermal relaxation \cite{moyses-v2, moyses-v3, tipler-v1, tipler-v2, callen-book}. Despite their distinct contexts, this processes are connected in a general way with the question of how systems evolve in time after a perturbation. In this work, we will focus on thermal relaxation processes where we have system interacting with an environment. According to thermodynamics, thermal relaxation is the exchange energy varying the system's temperature until it reaches a steady state. If the system is in a thermodynamic state near equilibrium, this evolution can be described using linear response theory, such as Onsager-Casemir and Kubo-Yokota-Nakajima laws \cite{lapolla-2020}. In this regime, the mean behaviour of statistical fluctuations is the same as the macroscopic system, being valid the assumption that the dynamics is memoryless and symmetric. Considering a thermal relaxation process, symmetry means that heating up or cooling down the system will take the same amount of time and follow the same thermodynamic trajectory \cite{ibanez-2024, lapolla-2020, seifert-2012, groot-2013}.

However, thermal relaxation symmetry can break for systems initially far from equilibrium and quantum systems, where linear response theory fails. One example of thermal relaxation asymmetry is the Mpemba effect. First observations of the Mpemba effect were that a hotter system (such as water) would cool down faster than a colder one \cite{mpemba-1969, kell-1969, teza-2026}. There is also the inverse Mpemba effect that occurs in heating processes instead of cooling \cite{teza-2026}. In general terms, Mpemba-like effects consist of a system further from equilibrium relaxing faster than another system closer to equilibrium. In this respect, Mpemba effects have been explored in the quantum regime \cite{krissia-2024, summer-2026, calabrese-2026, ares-2025, teza-2026}, such as the ergotropic Mpemba effect, where it was found that squeezed states relax faster than displaced states even when the ergotropy of squeezing is bigger (which means it is further from equilibrium) \cite{medina-2025}. Another relaxation asymmetry that emerges in the far-from-equilibrium regime is the heating-cooling asymmetry. Recent experiments using Brownian motion showed that heating is faster than cooling \cite{lapolla-2020, ibanez-2024, dieball-2023}. This was also observed to happen in quantum systems, such as a qubit, a quantum harmonic oscillator, and in quantum Brownian motion \cite{tejero-2025, manikandan-2021}. However, this behaviour is not universal, as it is possible to observe cooling relaxing faster than heating depending on the system's characteristics, for example, energy gap or degeneracy in systems with more than two levels \cite{vu-2021, bravetti-2025, teza-2026}. In addition, it is interesting that this heating-cooling asymmetry can also appear when comparing the usual Mpemba effect (cooling) with the inverse Mpemba effect (heating) \cite{teza-2026}.

A better understanding of relaxation asymmetries can be useful for state preparation \cite{verstraete-2009, harrington-2022}, optimal transport \cite{walker-2023}, heat engines \cite{liu-2024, liu-2026}, and other applications \cite{teza-2026}. For this purpose, beyond phenomenological characterisation, it is fundamental to study the physical mechanisms generating the asymmetries. To do this, previous studies developed an approach named thermal kinematics using the combination of stochastic thermodynamics and information geometry \cite{ito-2018, dechant-2020,ibanez-2024}. In thermal kinematics, thermodynamic states can be interpreted as points in a manifold, and we can define a distance between different states. As the system evolves in time, it will have a velocity and follow a thermodynamic trajectory in this manifold. Using this approach, Ref. \cite{ibanez-2024} proved that heating and cooling systems far from equilibrium follow distinct thermodynamic trajectories. Later, thermal kinematics was extended to the quantum regime, where this same result was obtained \cite{tejero-2025}. From a thermodynamic perspective, previous works showed that heating-cooling asymmetry can be explained by entropy production, in which heating is faster because it has a higher entropy production than cooling can dissipate heat \cite{tejero-misc-2025}. In addition to these works, understanding the consequences of quantumness in the heating-cooling asymmetry is still an open question. It is known that quantum correlations have consequences in thermodynamic processes, such as the case of anomalous heat exchange \cite{amikam-2020} and non-classical work extraction \cite{paternostro-2020, campaioli-2024}. Additionally, in the context of thermal relaxation, coherence in a qubit was shown to delay equilibration \cite{krissia-2024}. Therefore, the main question of this work is what quantumness do with the heating-cooling asymmetry.

There is a variety of systems and models we can use to explore this question. By observing different physical theories and phenomena, we see that discrete and continuous-variable models bring complementary explanations of nature. In addition, finding the analogous phenomena in both types of systems also enriches potential technological applications. Motivated by these considerations, this work explores the thermal relaxation asymmetry in both discrete and continuous-variable quantum systems. For the discrete case, we will focus on the thermal relaxation of two interacting qubits coupled to a global bath. Despite the variety of quantum properties we could investigate, we will focus only on quantum coherence and entanglement for this model. For the continuous-variable case, we focus on Gaussian states due to their relevance in quantum theory and experimental accessibility \cite{adesso-2016, stefanini-2026}.

In this work, we start discussing how open quantum systems evolve in time in Chap. \ref{cap:oqs}, where we derive the Lindblad master equation and present the thermalization dynamics called Davies map, that will be used in our results. Then, in Chap. \ref{cap:thermo}, we extend some thermodynamics concepts and quantities to the quantum regime, such as entropy and free energy. This chapter also introduces the thermal kinematics approach for both classical and quantum regime. Chap. \ref{cap:results-discrete} starts presenting the quantum properties that we will study for the discrete case. Next, we show our results for the heating-cooling asymmetry in two-interacting qubits model, where we investigate the effects of initial global coherences and correlations generated between the qubits during the dynamics. In Chap. \ref{cap:results-continuous}, we present the continuous-variable systems we will study that are Gaussian states. After showing the basic concepts, we discuss the thermalization of Gaussian states and reformulate the thermal kinematics approach for this case. With this, we show the results for the heating-cooling asymmetry in a thermal state, a displaced state, a squeezed state, and we show one way to invert the asymmetry. We end this work by summarising our results and suggesting future directions on this research topic.

